\section{Software Testing}

Software testing verifies that the AI Recipe Generator works according to the project requirements. The testing focused on authentication, protected access, recipe generation, pantry management, meal planning, shopping list creation, and profile updates.

\subsection{Testing Objectives}
\begin{itemize}
\item To verify that all implemented modules work according to the requirements.
\item To check user registration, login, and protected page access.
\item To validate recipe generation using user inputs.
\item To check database storage and retrieval of user data.
\item To ensure the application is user-friendly and responsive.
\end{itemize}

\subsection{Test Cases}
\begin{table}[h]
\centering
\caption{Test Cases}
\begin{tabularx}{\textwidth}{|p{1.5cm}|X|X|p{2cm}|}
\hline
\textbf{Test ID} & \textbf{Test Case} & \textbf{Expected Result} & \textbf{Status} \\
\hline
T1 & User registration with valid details. & Account should be created successfully. & Pass \\
\hline
T2 & User login with valid credentials. & User should log in and receive JWT token. & Pass \\
\hline
T3 & Login with invalid password. & Error message should be displayed. & Pass \\
\hline
T4 & Access protected page without token. & Access should be denied. & Pass \\
\hline
T5 & Generate recipe using ingredients. & Recipe should be generated successfully. & Pass \\
\hline
T6 & Save generated recipe. & Recipe should be saved in database. & Pass \\
\hline
T7 & Add pantry item. & Pantry item should be stored successfully. & Pass \\
\hline
T8 & Create meal plan. & Meal plan should be created successfully. & Pass \\
\hline
T9 & Generate shopping list. & Shopping list should be displayed correctly. & Pass \\
\hline
T10 & Update user profile. & Profile details should be updated successfully. & Pass \\
\hline
\end{tabularx}
\end{table}

\FloatBarrier
